\documentclass[lang=cn,color=blue,11pt]{elegantbook}
\usepackage{amsmath,amssymb,mathtools,booktabs}
\title{大学物理原创试题库}
\subtitle{题目选编}
\author{三千弱水 · Ruoshui Physics Bank}
\date{2026年10月}
\version{0.2}
\hypersetup{pdftitle={大学物理原创试题库 · 题目选编}}
\begin{document}
\maketitle
\tableofcontents
\mainmatter
\chapter{力学}

\noindent{\large\bfseries 【M007】 切向碰撞、扭转振子与粘滞阻尼}\quad{\small 难度 5/5 · 综合推导题}\par
\noindent\textbf{题目：}% M007: independent question
在光滑水平轨道上，质量为 $M$ 的滑块静止并与一端固定的轻弹簧（劲度系数为 $k$）连接，弹簧起初处于原长。质量为 $m$ 的小物块以速率 $v_0$ 沿弹簧轴线撞向滑块并粘在其上。设碰撞时间极短，碰撞期间弹簧冲量可忽略。求：(1) 碰撞后整体速度；(2) 弹簧最大压缩量；(3) 碰撞损失的机械能。

\par\smallskip\noindent\hrulefill\par\medskip
\noindent{\large\bfseries 【M008】 旋转圆环上珠子的分岔与能量交换}\quad{\small 难度 5/5 · 综合推导题}\par
\noindent\textbf{题目：}半径为 \(R\) 的竖直光滑圆环绕自身竖直直径以恒定角速度 \(\Omega\) 旋转。一颗质量为 \(m\) 的珠子无摩擦地在环上滑动，令 \(\theta\) 为珠子相对于最低点的有向角度（\(-\pi<\theta\le\pi\)）。重力加速度为 \(g\)。环由理想电机维持恒定 \(\Omega\)。
(1) 在共旋转系中建立有效势和运动方程；(2) 找出所有平衡位置及其稳定性，特别指出临界 \(\Omega\)；(3) 给出稳定非零平衡附近的小振动频率和两侧对称极小值之间的势垒；(4) 说明实验室系机械能为何通常不守恒，并求电机瞬时功率表达式。

\par\smallskip\noindent\hrulefill\par\medskip
\noindent{\large\bfseries 【M009】 反平方势修正下的近圆轨道与进动}\quad{\small 难度 5/5 · 综合推导题}\par
\noindent\textbf{题目：}% M009 · 题干
一根质量为 $m$、长度为 $L$ 的均匀细杆，可绕一端的固定光滑水平转轴在竖直平面内转动。细杆最初水平且静止，随后无初速释放。已知细杆关于端点的转动惯量 $I=mL^2/3$。忽略空气阻力，求细杆首次经过竖直向下位置瞬间的角速度大小及转轴对杆的竖直作用力，并说明其方向。

\par\smallskip\noindent\hrulefill\par\medskip
\noindent{\large\bfseries 【M010】 双圆轨道间的霍曼转移、冲量与能量预算}\quad{\small 难度 5/5 · 综合推导题}\par
\noindent\textbf{题目：}质量可忽略的航天器绕引力参数为 \(\mu=GM>0\) 的中心天体运动，初始处于半径为 \(r_1\) 的圆轨道，要转移到同一平面内半径 \(r_2>r_1\) 的圆轨道。忽略脉冲期间的质量损失与第三体作用，只允许在转移椭圆的近拱点和远拱点施加顺行瞬时切向速度增量。
(1) 从比机械能与角动量守恒求转移轨道半长轴以及两拱点速度；(2) 求两次脉冲 \(\Delta v_1,\Delta v_2\) 和总 \(\Delta v\)；(3) 求转移时间、两次脉冲导致的比轨道能增量，并检验总能量预算；(4) 令 \(r_2=r_1(1+\varepsilon)\) 且 \(\varepsilon\to0^+\)，求总 \(\Delta v\) 的首阶渐近式，并说明与圆轨道速度变化的关系。

\par\smallskip\noindent\hrulefill\par\medskip
\noindent{\large\bfseries 【M011】 有限振幅单摆的椭圆积分、周期修正与分界轨道}\quad{\small 难度 5/5 · 综合推导题}\par
\noindent\textbf{题目：}质量为 \(m\) 的质点固定在长度 \(\ell\) 的无质量刚性杆末端，杆可绕固定水平转轴在竖直平面内自由转动。质点从与最低点夹角 \(0<\theta_0<\pi\) 处静止释放，重力加速度为 \(g>0\)，忽略阻力。刚性杆可承受拉力和压力。
(1) 建立能量积分，给出最低点角速度；(2) 将振动周期写为第一类完全椭圆积分 \(K(k)\) 的形式，并明确参数 \(k\)；(3) 求 \(\theta_0\to0\) 时周期的首个非平凡修正，解释为何小振幅近似低估周期；(4) 求 \(\theta_0\to\pi^-\) 时周期的主导对数发散式，并求经过最低点瞬间杆对质点的径向作用力，说明刚性杆假设的用途。

\par\smallskip\noindent\hrulefill\par\medskip
\noindent{\large\bfseries 【M012】 二次空气阻力中的竖直抛体、逆运动与极值}\quad{\small 难度 5/5 · 综合推导题}\par
\noindent\textbf{题目：}% M012 · independent problem
在一维水平弹簧振子实验中，忽略摩擦，考虑弹簧与夹具的有效附加质量 $m_{\mathrm e}$，模型为
\[
T^2=\frac{4\pi^2}{k}(m+m_{\mathrm e}),
\]
其中 $m$ 为可更换砝码的质量。实验者以 SI 单位绘制 $T^2$ 对 $m$ 的图像，拟合直线为
\[
T^2=(0.80\,\mathrm{s^2/kg})m+0.16\,\mathrm{s^2}.
\]
求弹簧劲度系数 $k$ 和有效附加质量 $m_{\mathrm e}$，并写出所用的斜率、截距关系。

\par\smallskip\noindent\hrulefill\par\medskip
\noindent{\large\bfseries 【M013】 幂律中心势的圆轨道、稳定性及近心点进动}\quad{\small 难度 5/5 · 综合推导题}\par
\noindent\textbf{题目：}在平面内运动的质量 \(m>0\) 的质点受到吸引中心力 \(F_r=-k\,r^{p-1}\)，其中 \(k>0\)、实指数 \(p\in\mathbb R\)，角动量 \(L\ne0\)。当 \(p\ne0\) 时势能取 \(U(r)=kr^p/p\)；当 \(p=0\) 时定义 \(U(r)=k\ln(r/r_*)\)，其中 \(r_*>0\) 是任意参考长度。
(1) 写出径向有效势并由此求半径 \(r_0\) 的圆轨道所需角动量；(2) 严格判断圆轨道对小径向扰动稳定的 \(p\) 取值范围；(3) 在稳定范围内求径向振动与角向频率之比、相邻近心点间的方位角，判断对 \(p=-1\) 与 \(p=2\) 的闭合轨道结果；(4) 计算圆轨道的比动能及向心力平衡，说明为何 \(p\to0\) 的对数势极限是连续的。

\par\smallskip\noindent\hrulefill\par\medskip
\noindent{\large\bfseries 【M014】 水平匀加速支点下摆的有效重力、稳定平衡与张力}\quad{\small 难度 5/5 · 综合推导题}\par
\noindent\textbf{题目：}长度为 $\ell$ 的轻刚性杆悬挂质量为 $m$ 的小球，支点在竖直平面内以恒定水平加速度 $a>0$ 向 $+x$ 方向运动，重力为 $g>0$ 向下。摆角 $\theta$ 从竖直向下方向朝 $+x$ 测量。研究支点匀加速后的非惯性系运动。
(1) 建立摆的运动方程及守恒的有效机械能；(2) 求全部平衡角、稳定性以及稳定角附近小振动频率；(3) 支点原本静止、摆在 $\theta=0$ 处静止，突然开始加速度 $a$，求后续摆角振幅和两个转向角；(4) 求杆的瞬时径向作用力，特别核对初始时刻与稳定平衡时的值，并讨论 $a\to0$ 极限。

\par\smallskip\noindent\hrulefill\par\medskip
\noindent{\large\bfseries 【M015】 最速降线的变分方程、摆线参数与精确行程时间}\quad{\small 难度 5/5 · 综合推导题}\par
\noindent\textbf{题目：}% M015 · 题干
质量为 $m=0.020\,\mathrm{kg}$ 的小球水平射入质量为 $M=1.98\,\mathrm{kg}$ 的静止弹道摆并嵌入其中，摆块随后升高 $h=0.050\,\mathrm m$。设碰撞时间极短、悬线的水平冲量可忽略、上摆无耗散，取 $g=10\,\mathrm{m/s^2}$。求小球撞前速度，并说明两阶段分别适用何种守恒定律。

\par\smallskip\noindent\hrulefill\par\medskip
\noindent{\large\bfseries 【M016】 恒角速旋转杆上的径向滑珠与驱动功}\quad{\small 难度 5/5 · 综合推导题}\par
\noindent\textbf{题目：}% M016 · 题干
质量为 $m$ 的滑块接轻弹簧（劲度系数 $k$），在光滑水平面作简谐振动。以平衡位置为坐标原点，从 $x=A>0$ 处静止释放。求首次经过 $x=A/2$ 时的速度大小和此时动能与总机械能之比。

\par\smallskip\noindent\hrulefill\par\medskip
\noindent{\large\bfseries 【M017】 匀速排气火箭在重力场中的精确速度、高度与燃尽极限}\quad{\small 难度 5/5 · 综合推导题}\par
\noindent\textbf{题目：}% M017 · 题干
质点在固定半径 $r>0$ 的圆周上以恒定速率 $v>0$ 运动。在惯性参考系中，以下说法正确的是哪项？
\par\noindent\textbf{A.} 速率恒定，加速度恒为零
\par\noindent\textbf{B.} 加速度大小恒定，但方向不断改变
\par\noindent\textbf{C.} 加速度沿切线且与速度同向
\par\noindent\textbf{D.} 合外力始终沿速度方向

\par\smallskip\noindent\hrulefill\par\medskip
\chapter{电磁学}

\noindent{\large\bfseries 【E007】 滑动导体棒的自感反馈与能量耗散}\quad{\small 难度 5/5 · 综合推导题}\par
\noindent\textbf{题目：}% E007: independent question
平行板电容器的两极板尺寸为 $\ell\times w$，板间距 $d$，与恒压电源 $V$ 相连。一块厚度等于 $d$、相对介电常数为 $\kappa>1$ 的绝缘介质板沿宽度方向缓慢插入，插入长度为 $x$，其中 $0<x<w$。忽略边缘效应及摩擦。求：(1) 电容 $C(x)$；(2) 介质板受到的沿插入方向的静电力，并说明为什么不能直接以电容器储能的负梯度求恒压条件下的力。

\par\smallskip\noindent\hrulefill\par\medskip
\noindent{\large\bfseries 【E008】 同轴回流导线的磁场、磁能和内电感}\quad{\small 难度 4/5 · 综合推导题}\par
\noindent\textbf{题目：}% E008: independent question
质量为 $m$、电荷量为 $q>0$ 的粒子在真空中进入沿 $z$ 轴正方向的匀强磁场，磁感应强度大小为 $B>0$。初速度大小为 $v$，与磁场夹角为 $0<\alpha<\pi/2$。忽略辐射和重力，求螺旋轨道半径、回旋周期以及相邻两圈的轴向螺距。

\par\smallskip\noindent\hrulefill\par\medskip
\noindent{\large\bfseries 【E009】 正交电磁场中带黏滞阻力粒子的漂移}\quad{\small 难度 5/5 · 综合推导题}\par
\noindent\textbf{题目：}% E009 · 题干
一段长度为 $\ell$、电导率为常量 $\sigma>0$ 的均匀圆筒形导电介质填充在同轴圆柱面 $r=a$ 与 $r=b$ 之间，$0<a<b$。内外圆柱面被理想等势电极完全覆盖，电流 $I>0$ 沿径向向外稳恒流动，忽略两端边缘效应。求介质内电流密度、电场及两电极之间的等效电阻。

\par\smallskip\noindent\hrulefill\par\medskip
\noindent{\large\bfseries 【E010】 互感耦合 LC 回路的正规模、能量与拍动}\quad{\small 难度 5/5 · 综合推导题}\par
\noindent\textbf{题目：}% E010 · 题干
两只理想电容器的电容分别为 $C_1,C_2>0$，各自初始电压为 $V_1,V_2$（相对各自标记正端测量，允许不同），先彼此隔离。随后用具有微小但非零电阻的导线把同极性端子相连，待电流衰减至零。忽略辐射与电感储能，求末态共同电压和在导线中耗散的焦耳热。

\par\smallskip\noindent\hrulefill\par\medskip
\noindent{\large\bfseries 【E011】 径向非均匀自由电荷介质球的电势、束缚电荷与场能}\quad{\small 难度 5/5 · 综合推导题}\par
\noindent\textbf{题目：}% E011 · independent problem
真空中某任意形状的闭合高斯面内净电荷为 $+q$（$q>0$），高斯面外另有若干静止点电荷。设电通量方向按向外为正，求穿过整个闭合高斯面的总电通量。下列哪项正确？
\par\noindent\textbf{A.} 必为零，因为高斯面外存在电荷
\par\noindent\textbf{B.} 等于 $q/(4\pi\varepsilon_0)$
\par\noindent\textbf{C.} 等于 $q/\varepsilon_0$，与外部电荷的分布无关
\par\noindent\textbf{D.} 等于 $q/(\varepsilon_0 S)$，其中 $S$ 为高斯面的面积

\par\smallskip\noindent\hrulefill\par\medskip
\noindent{\large\bfseries 【E012】 幂律非均匀介电同轴电容器的场、电势与径向受力}\quad{\small 难度 5/5 · 综合推导题}\par
\noindent\textbf{题目：}一段长度 $\ell$ 的同轴圆柱电容器内外半径为 $0<a<b$，忽略端部效应。内导体总自由电荷 $+Q$，外导体总自由电荷 $-Q$，中间填充各向同性线性介质，径向介电常数按 $\epsilon(r)=\epsilon_a(r/a)^s$ 变化，其中 $\epsilon_a>0,\ s\in\mathbb R$。令内外导体电势差 $V=\phi(a)-\phi(b)$。
(1) 求介质内 $\boldsymbol D(r),\boldsymbol E(r)$ 及 $V$；(2) 求电容 $C(s)$，分别处理 $s=0$ 和 $s\ne0$；(3) 从 $u=\boldsymbol E\cdot\boldsymbol D/2$ 积分求场能并与 $Q^2/(2C)$ 核对；(4) 若端部不变、介质分布整体随 $s$ 被准静态调节而固定 $Q$，求 $U$ 对 $s$ 的变化率积分表达式，并讨论 $b/a\to1^+$ 的平板近似。

\par\smallskip\noindent\hrulefill\par\medskip
\noindent{\large\bfseries 【E013】 有限厚度电阻导体中的磁扩散、趋肤深度与功耗}\quad{\small 难度 5/5 · 综合推导题}\par
\noindent\textbf{题目：}% E013 · 独立题干
一矩形单匝闭合金属线框，沿运动方向的边长为 $d=0.30\,\mathrm m$、垂直运动方向的边长为 $\ell=0.20\,\mathrm m$，总电阻为 $R=0.50\,\Omega$。线框以恒定速率 $v=2.0\,\mathrm{m/s}$ 穿过一固定匀强磁场的直线边界，磁感应强度为 $B=0.50\,\mathrm T$，方向垂直线框平面。假设线框初始完全在场外、磁场区域足够宽、线框平面始终不变，并忽略自感及摩擦。仅讨论线框从开始进入磁场到刚好完全进入的阶段。求感应电动势大小、电流大小、磁场对线框的阻力大小以及该阶段焦耳热。

\par\smallskip\noindent\hrulefill\par\medskip
\noindent{\large\bfseries 【E014】 同轴传输线的 TEM 模、特征阻抗和负载反射}\quad{\small 难度 5/5 · 综合推导题}\par
\noindent\textbf{题目：}理想无损同轴线内外导体半径为 \(a,b\)（\(0<a<b\)），中间均匀介质介电常数为 \(\epsilon>0\)、磁导率为 \(\mu>0\)。忽略导体电阻及介质损耗，讨论 TEM 模沿 \(+z\) 传播。定义总电压 \(V\) 为内外导体电势差，电流 \(I\) 为内导体沿 \(+z\) 电流。
(1) 从横向静电和环向磁场求单位长度电容 \(C'\)、电感 \(L'\)；(2) 从电报方程得到传播速度和特征阻抗 \(Z_0\)；(3) 当末端接纯电阻负载 \(R_L>0\) 时求电压波反射系数及在开路、短路极限的相位；(4) 通过横截面积分坡印廷矢量证明瞬时电磁功率是 \(VI\)，并说明能量在介质中流动而不必“沿导体内部传播”。

\par\smallskip\noindent\hrulefill\par\medskip
\noindent{\large\bfseries 【E015】 并联 RLC 自由衰减的复特征根、储能与总耗散}\quad{\small 难度 5/5 · 综合推导题}\par
\noindent\textbf{题目：}% E015 · 题干
电容 $C>0$ 的理想电容器最初无电荷，经总电阻 $R>0$ 接至电压恒定的理想电源 $V>0$。忽略自感、辐射与漏电。求充电到稳态期间电源输出能量、最终电场能与电阻耗散热量，并讨论结果是否依赖电阻 $R$。

\par\smallskip\noindent\hrulefill\par\medskip
\noindent{\large\bfseries 【E016】 双层介质球形电容器的束缚电荷、电容与能量一致性}\quad{\small 难度 5/5 · 综合推导题}\par
\noindent\textbf{题目：}两个同心理想导体球面半径分别为 $a,b$（$0<a<b$），内球带自由电荷 $+Q$，外球带 $-Q$，$Q>0$。区域 $a<r<c$ 填充介电常数为 $\epsilon_1>\epsilon_0$ 的各向同性均匀线性电介质，区域 $c<r<b$ 填充介电常数为 $\epsilon_2>\epsilon_0$ 的另一介质，$a<c<b$；界面没有自由面电荷，介质均无自由体电荷。
(1) 由电位移高斯定律求两个介质区域的径向场及球面间电势差；(2) 求总电容与静电能，并通过场能积分独立检验；(3) 求 $r=a$、$r=c$ 和 $r=b$ 上所有介质束缚面电荷密度，验证合计的束缚电荷为零；(4) 检验 $\epsilon_1=\epsilon_2$ 的均匀介质极限以及 $c\to a^+$ 的极限，解释为何界面法向电场一般不连续。

\par\smallskip\noindent\hrulefill\par\medskip
\noindent{\large\bfseries 【E017】 直流阶跃驱动串联 RLC 的过冲、焦耳热与无阻尼奇异极限}\quad{\small 难度 5/5 · 综合推导题}\par
\noindent\textbf{题目：}% E017 · 题干
固定形状的一匝闭合线圈电阻 $R>0$ 且不变，忽略自感。若外磁通量随时间变化率的大小增至原来的两倍，则线圈的感应电动势和电流大小如何变化？
\par\noindent\textbf{A.} 感应电动势与电流大小均加倍
\par\noindent\textbf{B.} 感应电动势不变而电流加倍
\par\noindent\textbf{C.} 感应电流与磁通量变化率无关
\par\noindent\textbf{D.} 感应电流减为原来的一半

\par\smallskip\noindent\hrulefill\par\medskip
\chapter{热学}

\noindent{\large\bfseries 【T007】 导热可动活塞的不可逆终态与熵产生}\quad{\small 难度 5/5 · 综合推导题}\par
\noindent\textbf{题目：}一个刚性绝热容器总容积为 $V$，由可移动、导热、不透气的无热容活塞分隔为两室。室 $j$ 中有 $n_j$ mol 同种理想气体，初态温度 $T_j>0$、体积 $V_j>0$，其中 $V_1+V_2=V$。两室气体的摩尔定容热容为相同常数 $C_V$。活塞可由微弱耗散机制最终静止，容器及活塞不向外界传热或做功；初始和最终活塞均静止。
(1) 求最终共同温度和两室体积；(2) 求系统熵变 $\Delta S$ 并给出终态压力；(3) 从总熵极大原理独立推出热、力平衡条件；(4) 说明为何不能仅凭等温或等压过程公式描述整个不可逆演化。

\par\smallskip\noindent\hrulefill\par\medskip
\noindent{\large\bfseries 【T008】 有限热容热源驱动的最大可用功}\quad{\small 难度 5/5 · 综合推导题}\par
\noindent\textbf{题目：}% T008: independent question
一定量单原子理想气体按 $P$-$V$ 图上的闭合矩形作顺时针准静态循环：$A(V_0,P_0)\to B(V_0,2P_0)\to C(2V_0,2P_0)\to D(2V_0,P_0)\to A$。设各阶段可逆且仅存在体积功，求：(1) 一循环气体对外净做功；(2) 从外界吸收的总热量；(3) 热效率。

\par\smallskip\noindent\hrulefill\par\medskip
\noindent{\large\bfseries 【T009】 弹簧活塞的不可逆绝热平衡与熵极大}\quad{\small 难度 5/5 · 综合推导题}\par
\noindent\textbf{题目：}% T009 · 题干
某单原子理想气体经历准静态过程，满足 $PV^2=\text{常量}$。初态为 $(P_1,V_1,T_1)$，末态体积变为 $V_2=2V_1$。设气体物质的量不变，气体仅做体积功，取气体对外做功为正。求末态温度、该过程中气体对外做的功以及气体吸收的总热量，并判断其吸放热方向。

\par\smallskip\noindent\hrulefill\par\medskip
\noindent{\large\bfseries 【T010】 任意多方指数理想气体的热量符号与可逆绝热边界}\quad{\small 难度 4/5 · 综合推导题}\par
\noindent\textbf{题目：}% T010 · 题干
一台可逆制冷机工作于温度恒定的冷源 $T_c=270\,\mathrm K$ 与热源 $T_h=300\,\mathrm K$ 之间。某一循环从冷源吸热 $Q_c=900\,\mathrm J$。求制冷系数、外界输入的功以及向高温热源排放的热量。

\par\smallskip\noindent\hrulefill\par\medskip
\noindent{\large\bfseries 【T011】 范德瓦耳斯流体的临界点、旋节线与临界奇性}\quad{\small 难度 5/5 · 综合推导题}\par
\noindent\textbf{题目：}% T011 · independent problem
一绝热、刚性、密闭容器分为两部分，其中一侧有 $n$ 摩尔平衡态理想气体，初始体积为 $V$，另一侧为真空。移除隔板后气体自由膨胀至总体积 $2V$，达到新的平衡态。忽略气体与容器之间的能量交换。下列哪项关于气体末态温度和熵变的说法正确？
\par\noindent\textbf{A.} 温度下降，系统熵变为零
\par\noindent\textbf{B.} 温度不变，系统熵增加 $nRln2$
\par\noindent\textbf{C.} 温度上升，系统熵变为零
\par\noindent\textbf{D.} 温度不变，系统熵减少 $nRln2$

\par\smallskip\noindent\hrulefill\par\medskip
\noindent{\large\bfseries 【T012】 变热容有限热库驱动的可逆热机与最大功}\quad{\small 难度 5/5 · 综合推导题}\par
\noindent\textbf{题目：}% T012 · independent problem
在不与外界交换热量的理想热量计中，先装有 $0.20\,\mathrm{kg}$、$20.0^\circ\mathrm C$ 的冷水，热量计本身也处于该温度。随后迅速加入 $0.15\,\mathrm{kg}$、$80.0^\circ\mathrm C$ 的热水并充分搅拌，最终平衡温度为 $40.0^\circ\mathrm C$。两部分水的比热容均为 $c=4200\,\mathrm{J/(kg\cdot K)}$，忽略蒸发与相变。求热量计的总热容 $C_{\mathrm{cal}}$。

\par\smallskip\noindent\hrulefill\par\medskip
\noindent{\large\bfseries 【T013】 热容随温度线性增长的可逆绝热过程与熵函数}\quad{\small 难度 5/5 · 综合推导题}\par
\noindent\textbf{题目：}% T013 · 独立题干
一台循环热机只与温度保持恒定的两个热源交换热量，其高温热源温度 $T_h=600\,\mathrm K$，低温热源温度 $T_c=300\,\mathrm K$。热机每循环从高温热源吸热 $Q_h=800\,\mathrm J$。以下哪项正确描述热机每循环对外净做功的理论上限？
\par\noindent\textbf{A.} 最多 $200,mathrm J$
\par\noindent\textbf{B.} 最多 $400,mathrm J$
\par\noindent\textbf{C.} 最多 $800,mathrm J$
\par\noindent\textbf{D.} 只要循环足够缓慢，就能把吸热全部变为功

\par\smallskip\noindent\hrulefill\par\medskip
\noindent{\large\bfseries 【T014】 理想蒸汽近似下的克劳修斯—克拉佩龙曲线和熵跃变}\quad{\small 难度 5/5 · 综合推导题}\par
\noindent\textbf{题目：}% T014 · 独立题干
两段长度均为 $L$、横截面积均为 $A$ 的均匀细棒沿轴向紧密串联；左段材料导热系数为 $\kappa>0$，右段为 $2\kappa$。整体侧表面绝热，接触热阻可以忽略，两端分别维持在 $360\,\mathrm K$ 和 $300\,\mathrm K$。在稳态一维导热条件下，求两棒接合面的温度 $T_j$ 与通过两棒的热流功率 $\dot Q$。

\par\smallskip\noindent\hrulefill\par\medskip
\noindent{\large\bfseries 【T015】 范德瓦耳斯气体的低压焦耳—汤姆孙系数与反转温度}\quad{\small 难度 5/5 · 综合推导题}\par
\noindent\textbf{题目：}% T015 · 题干
某一理想气体温度由 $T$ 提升到 $4T$（均为开尔文温度），分子质量不变，服从经典麦克斯韦分布。气体分子的方均根速率变为原来的多少倍？
\par\noindent\textbf{A.} 变为四倍
\par\noindent\textbf{B.} 保持不变
\par\noindent\textbf{C.} 变为八倍
\par\noindent\textbf{D.} 变为两倍

\par\smallskip\noindent\hrulefill\par\medskip
\noindent{\large\bfseries 【T016】 自旋二能级顺磁体的配分函数、热容峰与绝热去磁}\quad{\small 难度 5/5 · 综合推导题}\par
\noindent\textbf{题目：}由 \(N\) 个互不相互作用的可区分自旋 \(1/2\) 磁矩组成理想顺磁体，每个自旋在恒定均匀正磁场 \(B>0\) 中只有两个能级 \(\varepsilon_\pm=\mp\mu B\)，其中 \(\mu>0\)。系统在热平衡温度 \(T>0\) 下服从正则分布，取玻尔兹曼常数 \(k_B\)，不考虑晶格热容、交换作用与量子简并。
(1) 求正则配分函数、平均内能与磁化强度；(2) 求定磁场热容与熵，分析高温和低温极限及热容为什么在两端均趋于零；(3) 求弱磁场磁化率并得到 Curie 定律；(4) 若通过理想可逆绝热过程将磁场从 \(B_i\) 缓慢降至 \(B_f\)（\(0<B_f<B_i\)），求温度变化关系，说明无法由该模型在有限正磁场下到达绝对零度。

\par\smallskip\noindent\hrulefill\par\medskip
\noindent{\large\bfseries 【T017】 黑体光子气体的辐射压力、绝热膨胀与熵约束}\quad{\small 难度 5/5 · 综合推导题}\par
\noindent\textbf{题目：}% T017 · 独立题干
$n$ 摩尔理想气体与温度为 $T>0$ 的恒温热源接触，经历缓慢可逆等温压缩，体积由 $V$ 降至 $V/3$。取气体对外做功为正，气体吸热为正。分别求该过程中气体的体积功 $W$、吸热量 $Q$ 与系统熵变 $\Delta S$，解释各量的符号。

\par\smallskip\noindent\hrulefill\par\medskip
\chapter{振动与波动}

\noindent{\large\bfseries 【W007】 双弦接点附加质量的波散射}\quad{\small 难度 5/5 · 综合推导题}\par
\noindent\textbf{题目：}% W007: independent question
在同一观测点，两列同频、同偏振方向且相干的单色光波单独存在时的光强分别为 $I_0$ 和 $4I_0$，其中 $I_0>0$。设两列波相位差为 $\delta$。求合成光强与 $\delta$ 的关系，以及可达到的最大和最小合成光强，并说明最小光强为何不为零。

\par\smallskip\noindent\hrulefill\par\medskip
\noindent{\large\bfseries 【W008】 耦合双振子的本征模与能量拍动}\quad{\small 难度 4/5 · 综合推导题}\par
\noindent\textbf{题目：}两相同质点质量为 \(m\)，各由一根刚度为 \(k>0\) 的弹簧连接到固定墙；两个质点之间另以刚度为 \(\kappa>0\) 的弹簧相连。设平衡位置两点位移为 \(x_1,x_2\)，所有运动均在同一水平方向且无摩擦。初态 \(x_1(0)=A\)、\(x_2(0)=0\)、\(\dot x_1(0)=\dot x_2(0)=0\)。
(1) 写出运动方程与两个正则频率；(2) 求两点位移全时解；(3) 在弱耦合 \(\kappa\ll k\) 时给出拍动的慢时间尺度，区分慢包络峰值与精确能量全转移；(4) 用总能量守恒和 \(\kappa\to0\) 检验答案。

\par\smallskip\noindent\hrulefill\par\medskip
\noindent{\large\bfseries 【W009】 三振子耦合本征模、能量分配与弱耦合拍动}\quad{\small 难度 5/5 · 综合推导题}\par
\noindent\textbf{题目：}% W009 · 题干
两根均匀柔软细弦在 $x=0$ 处平滑连接，张力相同为 $F>0$，左弦线密度为 $\mu_1$，右弦线密度为 $\mu_2=4\mu_1$。一列频率固定、振幅为 $A_i$ 的小振幅简谐横波自左向右正入射，接头质量可忽略且没有能量耗散。求位移波反射与透射的振幅系数 $A_r/A_i$、$A_t/A_i$，以及反射与透射平均功率占入射功率的比例。

\par\smallskip\noindent\hrulefill\par\medskip
\noindent{\large\bfseries 【W010】 受迫阻尼振子的幅值共振、功率共振与精确带宽}\quad{\small 难度 5/5 · 综合推导题}\par
\noindent\textbf{题目：}% W010 · 题干
某理想一维色散波的角频率与正波数满足 $\omega(k)=\alpha k^2$，其中 $\alpha>0$ 的量纲为 $\mathrm{m^2/s}$。一列窄带波包由集中在 $k_0>0$ 附近的波数组成。求其中心波数处的相速度与群速度，并讨论二者的比值及方向。

\par\smallskip\noindent\hrulefill\par\medskip
\noindent{\large\bfseries 【W011】 双原子弹簧链的色散关系、声学支和带隙}\quad{\small 难度 5/5 · 综合推导题}\par
\noindent\textbf{题目：}% W011 · independent problem
某一维简谐行波的位移函数为 $y(x,t)=A\cos(\omega t-kx)$，其中 $A,\omega,k$ 均为正数。关于波的传播方向和相速度，下列哪项正确？
\par\noindent\textbf{A.} 沿负 $x$ 方向，速率为 $omega/k$
\par\noindent\textbf{B.} 沿正 $x$ 方向，速率为 $omega/k$
\par\noindent\textbf{C.} 沿正 $x$ 方向，速率为 $k/omega$
\par\noindent\textbf{D.} 为驻波，其波节固定在空间中

\par\smallskip\noindent\hrulefill\par\medskip
\noindent{\large\bfseries 【W012】 带质量与弹簧耦合的弦端复反射系数}\quad{\small 难度 5/5 · 综合推导题}\par
\noindent\textbf{题目：}% W012 · independent problem
一根均匀细弦在恒定张力 $F$ 下形成稳定驻波。频率计测得驱动频率为 $f=200\,\mathrm{Hz}$，用刻度尺测出相邻两个驻波节点之间的间距为 $d=0.20\,\mathrm m$。弦的线密度为 $\mu=5.0\,\mathrm{g/m}$。假设弦满足 $v=\sqrt{F/\mu}$ 且频率与弦上行波频率相同。求弦中的张力。

\par\smallskip\noindent\hrulefill\par\medskip
\noindent{\large\bfseries 【W013】 端点有弹簧约束的弦本征频率与渐近边界条件}\quad{\small 难度 5/5 · 综合推导题}\par
\noindent\textbf{题目：}% W013 · 独立题干
长度为 $L$ 的细长空气柱，一端刚性封闭、另一端开口。考虑理想一维声波，且忽略开口端修正。关于其纵向声驻波的共振频率，下列哪一项正确？
\par\noindent\textbf{A.} 共振频率依次为基频的 $1,2,3$ 倍
\par\noindent\textbf{B.} 只有偶数倍基频可以共振
\par\noindent\textbf{C.} 前三个共振频率依次为基频的 $1,3,5$ 倍
\par\noindent\textbf{D.} 共振频率与空气柱长度无关

\par\smallskip\noindent\hrulefill\par\medskip
\noindent{\large\bfseries 【W014】 单原子弹簧链轻质量缺陷诱导的局域振动模}\quad{\small 难度 5/5 · 综合推导题}\par
\noindent\textbf{题目：}无限长单原子链中除 $n=0$ 处的质量为 $M$ 外，所有格点质量均为 $m>0$；每对相邻粒子由劲度系数为 $k>0$ 的相同轻弹簧连接，平衡间距为 $a>0$，忽略耗散。取缺陷比主体质量轻，即 $0<M<m$。
(1) 求无缺陷时的行波色散关系和声学频带上边缘；(2) 设缺陷诱导的局域振动形式为 $u_n(t)=A(-1)^n e^{-\lambda|n|}e^{-i\omega t}$，其中 $\lambda>0$，由非缺陷格点方程求 $\omega(\lambda)$；(3) 使用缺陷点方程解出 $\lambda$ 和 $\omega$ 的精确表达式，证明频率位于正常频带之外；(4) 研究 $M\to m^-$ 和 $M\to0^+$ 的局域长度与频率极限，解释重质量缺陷是否能以同样形式产生高频局域模。

\par\smallskip\noindent\hrulefill\par\medskip
\noindent{\large\bfseries 【W015】 张力随位置平方增长的非均匀弦精确本征谱}\quad{\small 难度 5/5 · 综合推导题}\par
\noindent\textbf{题目：}% W015 · 独立题干
某处声场两次测量的声强级分别为 $L_1=60\,\mathrm{dB}$、$L_2=80\,\mathrm{dB}$。参考声强相同，声强级定义为 $L=10\log_{10}(I/I_{\rm ref})\,\mathrm{dB}$。求两次实际声强的比值 $I_2/I_1$，并填入空白：$\displaystyle I_2/I_1=\underline{\hspace{1.4cm}}$。

\par\smallskip\noindent\hrulefill\par\medskip
\noindent{\large\bfseries 【W016】 有限异质弦段的双界面散射、透射共振与能流守恒}\quad{\small 难度 5/5 · 综合推导题}\par
\noindent\textbf{题目：}一根无限长、始终保持相同静张力 \(T>0\) 的无损细弦，其线密度在 \(x<0\) 与 \(x>d\) 两个半无限区段均为 \(\mu_1>0\)，而在 \(0<x<d\) 为 \(\mu_2>0\)，\(d>0\)。角频率 \(\omega>0\) 的单色横波从 \(x=-\infty\) 以时间因子 \(e^{-i\omega t}\) 入射，假定界面位移和由弦张力产生的横向力连续。
(1) 写出三段的波数、特征机械阻抗和单一界面的反射透射振幅；(2) 对有限中间区段的所有多次内反射求总反射幅度 \(r\) 与总透射幅度 \(t\)；(3) 证明无损结构的 \(|r|^2+|t|^2=1\)，求透射率与中间层相位的关系；(4) 找出完全透射的厚度/频率条件和最小透射率，并讨论 \(d\to0\)、\(\mu_1=\mu_2\) 极限。

\par\smallskip\noindent\hrulefill\par\medskip
\noindent{\large\bfseries 【W017】 固定边缘圆形膜的贝塞尔本征模、简并与能量正性}\quad{\small 难度 5/5 · 综合推导题}\par
\noindent\textbf{题目：}% W017 · 独立题干
两列相干、同频、同偏振方向的简谐波从两个同相波源出发到达同一观测点，传播介质相同，观测点处的波程差为 $\Delta r=3\lambda/2$。假定没有额外反射相变，哪项描述正确？
\par\noindent\textbf{A.} 总是相长干涉
\par\noindent\textbf{B.} 程差与干涉强弱无关
\par\noindent\textbf{C.} 相消干涉，若等振幅可完全抵消
\par\noindent\textbf{D.} 该程差对应相位差为 $2pi$

\par\smallskip\noindent\hrulefill\par\medskip
\chapter{光学}

\noindent{\large\bfseries 【O007】 分离薄透镜系统的矩阵成像与无焦条件}\quad{\small 难度 5/5 · 综合推导题}\par
\noindent\textbf{题目：}% O007: independent question
强度为 $I_0$ 的自然光依次通过三片理想线偏振片。三片偏振轴相对于第一片分别为 $0^\circ$、$30^\circ$ 和 $90^\circ$。忽略反射和吸收以外的偏振损耗。求最终透射光强；若取走中间的偏振片，最终透射光强又是多少？

\par\smallskip\noindent\hrulefill\par\medskip
\noindent{\large\bfseries 【O008】 双缝干涉的单缝包络、缺级与相移}\quad{\small 难度 4/5 · 综合推导题}\par
\noindent\textbf{题目：}% O008: independent question
迈克耳孙干涉仪在空气中使用真空波长近似为 $\lambda=600\,\mathrm{nm}$ 的单色光。缓慢沿光束方向移动其中一面反射镜，位移大小为 $\Delta d=0.150\,\mathrm{mm}$。假设空气折射率可取1、干涉始终可观测，问某一固定视场位置共经过多少条明暗周期（完整干涉条纹）？

\par\smallskip\noindent\hrulefill\par\medskip
\noindent{\large\bfseries 【O009】 法布里—珀罗腔的线型、自由光谱区与半高宽}\quad{\small 难度 4/5 · 综合推导题}\par
\noindent\textbf{题目：}% O009 · 题干
一块透射光栅的相邻刻线间距为 $d=1.50\,\mathrm{\mu m}$，使用波长 $\lambda=600\,\mathrm{nm}$ 的平行单色光垂直入射。忽略单缝衍射包络导致的缺级，求在传播方向可观察到的最大衍射级次的绝对值，以及第二级亮纹相对光栅法线的衍射角。

\par\smallskip\noindent\hrulefill\par\medskip
\noindent{\large\bfseries 【O010】 高斯光束通过薄透镜后的复束参数与束腰位置}\quad{\small 难度 5/5 · 综合推导题}\par
\noindent\textbf{题目：}在均匀无耗散介质中传播的轴对称单模近轴高斯光束，其真空波长为 $\lambda_0>0$、介质折射率为 $n>0$。原束腰半径 $w_0$ 位于 $z=0$，瑞利长度 $z_R=\pi n w_0^2/\lambda_0$。一理想薄凸透镜焦距 $f>0$ 位于 $z=d>0$，两侧介质相同，忽略像差、衍射外的损耗。
(1) 写出透镜前复束参数 $q_{\rm in}$、束宽 $w(z)$ 和波前曲率 $R(z)$；(2) 用自由传播与薄透镜 ABCD 变换求紧靠透镜后的 $q_{\rm out}$，分离实部和虚部；(3) 求新束腰相对透镜的位置 $z_w$、新瑞利长度 $z_R'$ 与束腰半径 $w_0'$，判断何时新束腰位于透镜右侧；(4) 讨论 $z_R\gg\max(d,f)$ 的近似准直输入极限，并核对维度与焦距标度。

\par\smallskip\noindent\hrulefill\par\medskip
\noindent{\large\bfseries 【O011】 有限谱宽迈克耳孙干涉的相干包络与条纹可见度}\quad{\small 难度 4/5 · 综合推导题}\par
\noindent\textbf{题目：}一理想迈克耳孙干涉仪在空气中工作，光源为中心频率 \(\nu_0>0\) 的近似准单色非偏振光。假定各频率分量互不相干，光强谱在 \([\nu_0-\Delta\nu/2,\nu_0+\Delta\nu/2]\) 上均匀分布，区间外为零，且 \(0<\Delta\nu\ll\nu_0\)。两臂返回到同一探测点时，单独打开任一臂的积分光强相同为 \(I_0\)。定义两臂之间光程差 \(\ell\)、时间差 \(\tau=\ell/c\)。
(1) 对频率积分，严格求总探测光强 \(I(\tau)\)；(2) 求条纹可见度、第一次完全失去可见度时的光程差，以及为何之后仍可能出现旁瓣；(3) 当移动一面反射镜距离 \(x\) 时求 \(\ell\) 的变化，得到中心条纹周期与第一个可见度零点对应的镜面位移；(4) 以 \(\lambda_0=c/\nu_0\) 和窄带谱宽 \(\Delta\lambda\) 表示一阶相干长度，说明近似的条件。

\par\smallskip\noindent\hrulefill\par\medskip
\noindent{\large\bfseries 【O012】 有间距双透镜系统的主平面、后焦点与望远镜条件}\quad{\small 难度 5/5 · 综合推导题}\par
\noindent\textbf{题目：}% O012 · independent problem
单色平行光垂直照射一个宽度为 $a$ 的单缝，在缝后 $D=2.00\,\mathrm m$ 处的观察屏上，测得中央明纹中心到第一个暗纹中心的距离为 $x_1=5.00\,\mathrm{mm}$。光波长为 $\lambda=500\,\mathrm{nm}$。在近轴条件 $x_1\ll D$ 下，估计单缝缝宽 $a$，并说明所用近似。

\par\smallskip\noindent\hrulefill\par\medskip
\noindent{\large\bfseries 【O013】 二阶群速度色散下高斯脉冲的展宽与线性啁啾}\quad{\small 难度 5/5 · 综合推导题}\par
\noindent\textbf{题目：}% O013 · 独立题干
一理想平面透射光栅被单色准直光均匀照亮，共有 $N=1500$ 条刻线参与衍射。在中心波长约为 $\lambda=600\,\mathrm{nm}$ 的二级主极大方向，假设二级光谱存在且仪器其他部件不限制分辨率。按瑞利判据，求此处两条接近谱线所需的最小波长差。

\par\smallskip\noindent\hrulefill\par\medskip
\noindent{\large\bfseries 【O014】 有限数目等宽狭缝光栅的衍射包络、缺级与分辨本领}\quad{\small 难度 5/5 · 综合推导题}\par
\noindent\textbf{题目：}理想平面透射光栅由 $N\ge2$ 个完全相同的矩形狭缝组成，每缝宽 $a>0$、相邻缝中心距 $d>a$，且均匀相干单色平面波波长为 $\lambda>0$ 沿法线入射。只研究标量夫琅禾费衍射，在远场方向 $\theta$ 观察。
(1) 从单缝振幅积分和 $N$ 个等距相干波的几何级数求归一化光强函数；(2) 给出光栅主极大级次条件和单缝包络零点，严格推导缺级条件；(3) 在 $N\gg1$、第 $m\ne0$ 级不缺级且 $\cos\theta_m\ne0$ 时估算主峰第一个零点与主峰角宽度；(4) 使用瑞利判据推导光栅分辨本领 $\lambda/\Delta\lambda\simeq |m|N$，说明此式的适用条件与狭缝宽度影响。

\par\smallskip\noindent\hrulefill\par\medskip
\noindent{\large\bfseries 【O015】 四分之一波片对椭圆偏振的转换及分析器旋转测量}\quad{\small 难度 5/5 · 综合推导题}\par
\noindent\textbf{题目：}一束完全偏振的单色平面光以单位归一化 Jones 振幅 \(\boldsymbol E_{\rm in}=(\cos\theta,\sin\theta)^T\) 入射一理想四分之一波片，\(x\) 轴是快轴、\(y\) 轴是慢轴。以复时间因子 \(e^{-i\omega t}\) 为约定，波片的 Jones 矩阵取 \(J=\operatorname{diag}(1,i)\)。输出光随后通过透射轴与 \(x\) 方向夹角为 \(\alpha\) 的理想线偏振分析器。忽略全部非理想损耗。
(1) 求波片输出 Jones 向量及两分量的相位差；(2) 推导任意 \(\theta\) 下偏振椭圆的轴向投影幅度，给出输出为圆偏振和线偏振的精确条件；(3) 计算分析器透射光强占波片前入射强度的比例，给出其对 \(\alpha\) 的规律；(4) 当 \(\theta=\pi/4\) 时证明分析器无论怎样旋转都只透射一半强度，并解释为何这并不表示输出光是不偏振光。

\par\smallskip\noindent\hrulefill\par\medskip
\noindent{\large\bfseries 【O016】 单层增透膜的多界面复反射、四分之一波厚度与带宽}\quad{\small 难度 5/5 · 综合推导题}\par
\noindent\textbf{题目：}% O016 · 独立题干
迈克耳孙干涉仪在空气中以一束近似单色激光工作，空气折射率取1。将一面反射镜沿光束方向缓慢移动 $\Delta d=0.0750\,\mathrm{mm}$，某固定视场位置记录到 $N=250$ 次完整明暗周期通过。忽略振动误差，求激光波长；说明为什么镜面移动量要乘以2。

\par\smallskip\noindent\hrulefill\par\medskip
\noindent{\large\bfseries 【O017】 色散棱镜的最小偏向角、对称驻值与光谱角色散}\quad{\small 难度 5/5 · 综合推导题}\par
\noindent\textbf{题目：}一块顶角为 \(A\)（\(0<A<\pi\)）的理想均匀三棱镜处于空气中，折射率为 \(n(\lambda)>1\)，忽略反射损耗和吸收，仅研究从一个侧面入射并从另一侧面出射、无全反射的几何光线。入射角 \(i\)、出射角 \(e\) 均由对应表面法线量起；棱镜内两个折射角为 \(r_1,r_2\)，满足 \(r_1+r_2=A\)。
(1) 建立偏向角 \(D=i+e-A\) 与 Snell 关系；(2) 严格证明在满足 \(n\sin(A/2)<1\) 的对称传播分支上，偏向角在 \(r_1=r_2=A/2\) 取得局部极小，求 \(D_{\min}(\lambda)\)；(3) 求最小偏向角随波长的导数，并对 \(n(\lambda)=n_\infty+C/\lambda^2\)（\(C>0\)）判断角色散方向；(4) 求 \(A\to0^+\) 时的薄棱镜近似以及对称点附近偏向角的二阶展开，说明结果的几何适用条件。

\par\smallskip\noindent\hrulefill\par\medskip
\chapter{近代物理}

\noindent{\large\bfseries 【Q007】 相对论两体衰变的能量区间与阈值}\quad{\small 难度 5/5 · 综合推导题}\par
\noindent\textbf{题目：}静质量为 $M$ 的不稳定粒子衰变为静质量 $m_1,m_2$ 的两个粒子，假设可视为真空中的自由两体衰变，所有质量非负且 $M>0$。在母粒子静止系记子粒子 $1$ 的能量、动量大小为 $E_1^*,p^*$；在实验室系母粒子沿 $+x$ 以 $V$ 匀速运动（$0\le V<c$），令 $\theta^*$ 为子粒子 $1$ 在母粒子静止系内相对于 $+x$ 的发射角。
(1) 严格给出衰变允许的阈值条件；(2) 求 $E_1^*,E_2^*,p^*$；(3) 求实验室系的 $E_1(\theta^*)$、$p_{1x}(\theta^*)$ 及能量上下界；(4) 检查阈值 $M=m_1+m_2$ 和静止母粒子 $V=0$ 两个极限。

\par\smallskip\noindent\hrulefill\par\medskip
\noindent{\large\bfseries 【Q008】 无限深势阱突然扩张后的量子投影}\quad{\small 难度 5/5 · 综合推导题}\par
\noindent\textbf{题目：}一维无限深势阱初始区间为 $0<x<L$，质量为 $m$ 的粒子最初处于基态 $\psi(x,0^-)=\sqrt{2/L}\sin(\pi x/L)$。在 $t=0$ 突然将右墙移至 $x=2L$，随后新势阱保持不变。移动足够快，可使用突然近似：$\psi(x,0^+)=\psi(x,0^-)$（并在 $L<x<2L$ 延拓为 $0$）。
(1) 求新阱的本征能量与初态在各新本征态上的展开系数 $c_n$；(2) 求测得新阱能级 $n=2$ 的概率（$n=2$ 是第一激发态）和所有偶数 $n\ne2$ 的概率；(3) 写出 $t>0$ 波函数并计算骤变后能量期望，解释与旧能量的关系；(4) 判断新哈密顿量能量方差是否有限，并说明原因。

\par\smallskip\noindent\hrulefill\par\medskip
\noindent{\large\bfseries 【Q009】 自旋测量、拉莫尔进动与非选择性退相干}\quad{\small 难度 4/5 · 综合推导题}\par
\noindent\textbf{题目：}一个自旋 \(1/2\) 粒子初态为 \(|+z\rangle\)。先测量沿单位方向 \(\boldsymbol n=(\sin\theta,0,\cos\theta)\)（\(0<\theta<\pi\)）的自旋投影 \(\boldsymbol n\cdot\boldsymbol\sigma\)，其本征值为 \(\pm1\)。此后粒子在均匀磁场中按哈密顿量 \(H=-\hbar\omega\sigma_z/2\) 演化时间 \(t\ge0\)，\(\omega>0\)，再测量 \(\sigma_x\)。
(1) 求首次测到 \(+1\) 的概率及相应归一化后测态；(2) 在选择首次 \(+1\) 结果的子样本中，求末次测到 \(\sigma_x=+1\) 的条件概率；(3) 若首次测量仍实施但不记录结果，利用密度矩阵求末次测量概率，并与未经首次测量的情况比较；(4) 对 \(\theta=\pi/2\) 给出三个概率随时间的差异，说明为何丢弃测量记录不能等价于未作测量。

\par\smallskip\noindent\hrulefill\par\medskip
\noindent{\large\bfseries 【Q010】 恒定固有加速度运动的世界线、快度与双时钟}\quad{\small 难度 5/5 · 综合推导题}\par
\noindent\textbf{题目：}% Q010 · 题干
一质量为 $m$ 的粒子相对于某惯性实验室系以速率 $v=3c/5$ 做直线运动。求其洛伦兹因子 $\gamma$、动能 $K$ 和相对论动量大小 $p$，并用结果指出为什么不能套用非相对论动量公式 $p=mv$。

\par\smallskip\noindent\hrulefill\par\medskip
\noindent{\large\bfseries 【Q011】 有失谐双能级系统的拉比振荡与量子态保真度}\quad{\small 难度 5/5 · 综合推导题}\par
\noindent\textbf{题目：}% Q011 · independent problem
某金属在固定频率 $f>f_0$ 的单色光照射下发生普通单光子光电效应。保持金属表面和入射频率不变，增大照射光强，并假设收集条件和量子效率不变、未发生空间电荷限制。下列哪项正确？
\par\noindent\textbf{A.} 最大初动能增大，遏止电压增大
\par\noindent\textbf{B.} 饱和光电流增大，但遏止电压不变
\par\noindent\textbf{C.} 逸出功随入射光强增大
\par\noindent\textbf{D.} 只要光强足够大，低于阈频的光仍必然产生单光子光电效应

\par\smallskip\noindent\hrulefill\par\medskip
\noindent{\large\bfseries 【Q012】 量子势阶跃的传播反射、穿透衰减与概率流}\quad{\small 难度 5/5 · 综合推导题}\par
\noindent\textbf{题目：}% Q012 · independent problem
一放射性样品的计数率在两次测量中分别为 $850\,\mathrm{s^{-1}}$ 和 $250\,\mathrm{s^{-1}}$，两次测量相隔 $1.00\,\mathrm h$。同时独立测得环境本底计数率始终为 $50\,\mathrm{s^{-1}}$。假设仪器死时间效应可忽略、检测效率恒定、样品只有一种指数衰变核素，求核素的半衰期。

\par\smallskip\noindent\hrulefill\par\medskip
\noindent{\large\bfseries 【Q013】 量子自旋回波、相位反演与非均匀失谐退相干}\quad{\small 难度 5/5 · 综合推导题}\par
\noindent\textbf{题目：}理想自旋 \(1/2\) 粒子群起初均处于 \(|+x\rangle\) 态。每个粒子在相同旋转坐标系中受到纯静态失谐哈密顿量 \(H_\delta=\hbar\delta\sigma_z/2\)，其中每个粒子的 \(\delta\) 固定不变，但整个粒子群的失谐服从均值零、标准差 \(\sigma>0\) 的高斯分布。忽略自旋间相互作用以及随机时间涨落。
(1) 对给定 \(\delta\) 求自由演化时间 \(t\) 后 \(\langle\sigma_x\rangle_\delta\)；(2) 对高斯分布求群平均横向相干信号，得出非均匀退相干时间标度；(3) 在自由演化 \(\tau\) 后施加理想瞬时 \(\pi\) 脉冲 \(U_\pi=-i\sigma_x\)，再自由演化 \(\tau\)，求单粒子的总算符并证明 \(2\tau\) 时完全回波；(4) 解释为何自由感应衰减可由回波恢复，并讨论如果每个粒子的失谐在两个时间段随机且彼此独立时，理想回波是否仍必然完全恢复。

\par\smallskip\noindent\hrulefill\par\medskip
\noindent{\large\bfseries 【Q014】 谐振子频率骤变后的激发数、基态重叠与压缩振荡}\quad{\small 难度 5/5 · 综合推导题}\par
\noindent\textbf{题目：}% Q014 · 独立题干
飞船相对于实验室沿 $+x$ 方向匀速飞行，速率为 $v=\frac12c$。飞船中的装置沿同一方向向前发射一粒子，粒子相对于飞船的速率为 $u'=\frac35c$。忽略引力，求实验室测得粒子的速率，并将结果化为 $c$ 的有理数倍。

\par\smallskip\noindent\hrulefill\par\medskip
\noindent{\large\bfseries 【Q015】 相对论完全非弹性碰撞的复合静质量与内部能}\quad{\small 难度 5/5 · 综合推导题}\par
\noindent\textbf{题目：}% Q015 · 独立题干
某金属的光电效应实验中，用入射频率 $f_1=6.0\times10^{14}\,\mathrm{Hz}$ 测得遏止电压 $U_1=0.50\,\mathrm V$；用 $f_2=8.0\times10^{14}\,\mathrm{Hz}$ 测得 $U_2=1.30\,\mathrm V$。假设满足 $eU_s=hf-hf_0$，忽略接触电势修正。求该实验拟合得到的 $h/e$（以 $\mathrm{V\cdot s}$ 为单位）及金属阈频 $f_0$。

\par\smallskip\noindent\hrulefill\par\medskip
\noindent{\large\bfseries 【Q016】 量子谐振子相干态的泊松统计、经典轨迹与最小不确定性}\quad{\small 难度 5/5 · 综合推导题}\par
\noindent\textbf{题目：}% Q016 · 质量复核修订题干
对一非相对论量子粒子，一维位置和动量的标准差满足 $\Delta x\,\Delta p\geq\hbar/2$。若实验设计将允许的位置标准差从 $a>0$ 放宽到 $2a$，比较在两种给定位置标准差条件下**理论最小动量标准差的下界**，新下界相当于原下界的多少倍？
\par\noindent\textbf{A.} 动量不确定度下界变为原来的四倍
\par\noindent\textbf{B.} 动量不确定度下界变为原来的两倍
\par\noindent\textbf{C.} 动量不确定度下界变为原来的一半
\par\noindent\textbf{D.} 动量不确定度下界不变

\par\smallskip\noindent\hrulefill\par\medskip
\noindent{\large\bfseries 【Q017】 一维有限方势阱的偶奇束缚态、阈值数目与弱束缚渐近}\quad{\small 难度 5/5 · 综合推导题}\par
\noindent\textbf{题目：}质量为 $m>0$ 的非相对论粒子处在对称有限方势阱
\[
V(x)=\begin{cases}-V_0,&|x|<a,\\0,&|x|\ge a,\end{cases}\qquad V_0>0,\ a>0 .
\]
只研究平方可积束缚态 $-V_0<E<0$，忽略自旋与其他作用。
(1) 利用中心对称宇称和 $x=\pm a$ 处波函数及导数连续性，建立偶态、奇态的能级超越方程；(2) 用无量纲阱强度 $\lambda=a\sqrt{2mV_0}/\hbar$ 求束缚态数目，并说明阈值处为何不把 $E=0$ 计为束缚态；(3) 给出偶、奇态波函数的归一化常数表达式，证明能级均简单而没有奇偶简并；(4) 在 $V_0\to\infty$ 的固定低能级序号极限及 $V_0\to0^+$ 的浅阱极限求能量的主导近似。

\par\smallskip\noindent\hrulefill\par\medskip
\end{document}
